\begin{recipe}{Bananenijs}
\kicker[Dessert,Vegetarisch]{Dessert · vegetarisch}
\index[register]{Bananenijs}
\index[register]{Banaan}

\meta{5 minuten + een dag invriezen \dvd Voor 2 personen}

\lettrine{H}{eb} je nog een paar bananen liggen die net iets te bruin zijn om zo op
te eten? Snijd ze in plakjes, vries ze in en draai er met de blender in een paar
minuten een romig toetje van. Handig om overrijpe bananen niet te hoeven weggooien,
en kinderen vinden het meestal prima.

\blockrule{Ingrediënten}
\begin{ingredients}
  \ing{2}{rijpe bananen}\index[register]{Banaan}
  \ing{1 scheutje}{melk}
  \ingb{halve el cacaopoeder, optioneel}
\end{ingredients}

\blockrule{Bereiding}
\tip{Zet een zakje geschilde, in plakjes gesneden banaan standaard in de vriezer
zodra je bananen te bruin beginnen te worden. Dan heb je altijd voorraad voor dit
ijs.}
\begin{steps}
  \item Schil de bananen en snijd ze in plakjes. Doe ze in een ziplock zak en vries
        ze minstens een dag in.
  \item Haal de bevroren bananenplakjes uit de vriezer en doe ze in de blender met
        een scheutje melk.
  \item Voeg naar smaak de cacaopoeder toe en mix tot een gladde, romige massa die
        aanvoelt als roomijs. Voeg eventueel nog een scheutje melk toe als het
        moeilijk mixt.
  \item Serveer meteen.
\end{steps}
\end{recipe}
